\begin{tomtat}
	\begin{minipage}[t]{0.3\textwidth}
		\vspace{0pt}
		\begin{tomtat}
			\begin{center}
				\textbf{ĐỀ 28}
			\end{center}
			
		\end{tomtat}
	\end{minipage}
	\hfill
	\begin{minipage}[t]{0.5\textwidth}
		\vspace{0pt}
		\begin{tomtat}
			\begin{center}
				\includegraphics[scale=0.08]{images/28}
			\end{center}
		\end{tomtat}
	\end{minipage}
\end{tomtat}
\textbf{I. PHẦN TRẮC NGHIỆM}
\begin{ex}
	Phương trình nào sau đây là phương trình bậc nhất hai ẩn?
	\choice{$0x-3=0$}{$5x-2y=0$}{$x^2+2x-1=0$}{$x^2+1=0$}
	\loigiai{
		Phương trình bậc nhất hai ẩn có dạng $ax+by=c$, trong đó a, b, c là các số thực, và a, b không đồng thời bằng 0.
		A. $0x-3=0 \Leftrightarrow -3=0$ (Vô lý). Đây không phải là phương trình bậc nhất hai ẩn.
		B. $5x-2y=0$. Đây là phương trình bậc nhất hai ẩn với $a=5, b=-2, c=0$.
		C. $x^2+2x-1=0$. Đây là phương trình bậc hai một ẩn.
		D. $x^2+1=0$. Đây là phương trình bậc hai một ẩn (vô nghiệm trên tập số thực).
		Vậy phương trình $5x-2y=0$ là phương trình bậc nhất hai ẩn.
		\\ \textbf{Đáp Án B}
	}
\end{ex}

\begin{ex}
	Tập nghiệm của bất phương trình: $x-2>0$ là:
	\choice{$\{x \in \mathbb{R} / x > 2\}$}{$\{x \in \mathbb{R} / x < 2\}$}{$\{x \in \mathbb{R} / x \le 2\}$}{$\{x \in \mathbb{R} / x \ge -2\}$}
	\loigiai{
		Ta có: $x-2 > 0$
		$\Leftrightarrow x > 2$.
		Vậy tập nghiệm của bất phương trình là $\{x \in \mathbb{R} / x > 2\}$.
		\\ \textbf{Đáp Án A}
	}
\end{ex}

\begin{ex}
	Giá trị của biểu thức $A=5\sqrt{2}-\sqrt{8}$ bằng:
	\choice{$\sqrt{5}$}{$2$}{$3\sqrt{2}$}{$\sqrt{2}$}
	\loigiai{
		Ta có:
		\begin{align*} A &= 5\sqrt{2} - \sqrt{8} \\ &= 5\sqrt{2} - \sqrt{4 \cdot 2} \\ &= 5\sqrt{2} - 2\sqrt{2} \\ &= (5-2)\sqrt{2} \\ &= 3\sqrt{2} \end{align*}
		\\ \textbf{Đáp Án C}
	}
\end{ex}

\begin{ex}
	Cho tam giác ABC vuông cân tại A, khi đó giá trị lượng giác $\tan ABC$ là:
	\choice{$-\frac{\sqrt{2}}{2}$}{$\frac{1}{2}$}{$1$}{$\frac{\sqrt{2}}{2}$}
	\loigiai{
		Tam giác ABC vuông cân tại A nên $\widehat{B} = \widehat{ABC} = 45^\circ$.
		Do đó, $\tan ABC = \tan B = \tan 45^\circ = 1$.
		\\ \textbf{Đáp Án C}
	}
\end{ex}

\begin{ex}
	Độ dài cung nhỏ AB của đường tròn (O; 2cm) biết số đo $\widehat{AOB}=60^\circ$ là:
	\choice{$3\pi$}{$10\pi$}{$\frac{2\pi}{3}$}{$\frac{\pi}{3}$}
	\loigiai{
		Độ dài cung tròn được tính bằng công thức $l = \frac{\pi R n}{180}$, trong đó R là bán kính và n là số đo góc ở tâm (tính bằng độ).
		Với $R=2$ cm và $n=60^\circ$, ta có:
		$l = \frac{\pi \cdot 2 \cdot 60}{180} = \frac{120\pi}{180} = \frac{2\pi}{3}$ (cm).
		\\ \textbf{Đáp Án C}
	}
\end{ex}

\begin{ex}
	Cho hình nón có bán kính đáy $r=6$ cm và chiều cao $h=8$ cm. Diện tích xung quanh của hình nón là:
	\choice{$48\pi \text{ cm}^2$}{$60\pi \text{ cm}^2$}{$96\pi \text{ cm}^2$}{$120\pi \text{ cm}^2$}
	\loigiai{
		Đầu tiên, ta cần tính độ dài đường sinh $l$ của hình nón.
		Áp dụng định lý Pitago trong tam giác vuông tạo bởi chiều cao, bán kính đáy và đường sinh:
		$l^2 = r^2 + h^2 = 6^2 + 8^2 = 36 + 64 = 100$.
		$l = \sqrt{100} = 10$ (cm).
		Diện tích xung quanh của hình nón được tính bằng công thức $S_{xq} = \pi r l$.
		$S_{xq} = \pi \cdot 6 \cdot 10 = 60\pi \text{ (cm}^2\text{)}$.
		\\ \textbf{Đáp Án B}
	}
\end{ex}
\begin{ex}
	Trong các phương trình sau, phương trình nào là phương trình bậc hai một ẩn?
	\choice{$2x-3y=0$}{$x^2-\sqrt{x}+4=0$}{$x^2+2x+1=0$}{$0x^2+2x-4=0$}
	\loigiai{
		Phương trình bậc hai một ẩn có dạng $ax^2+bx+c=0$ với $a \neq 0$.
		A. $2x-3y=0$ là phương trình bậc nhất hai ẩn.
		B. $x^2-\sqrt{x}+4=0$ không phải là phương trình bậc hai vì chứa số hạng $\sqrt{x}$.
		C. $x^2+2x+1=0$ là phương trình bậc hai một ẩn (với $a=1 \neq 0, b=2, c=1$).
		D. $0x^2+2x-4=0 \Leftrightarrow 2x-4=0$, đây là phương trình bậc nhất một ẩn.
		Vậy phương trình $x^2+2x+1=0$ là phương trình bậc hai một ẩn.
		\\ \textbf{Đáp Án C}
	}
\end{ex}

\begin{ex}
	Một hộp chứa 4 quả cầu cùng loại trong đó có 1 quả cầu đỏ, 1 quả cầu xanh và 2 quả cầu vàng. Chọn ngẫu nhiên đồng thời ra hai quả cầu. Xác suất của biến cố “Chọn được 1 quả cầu đỏ và 1 quả cầu vàng” là:
	\choice{$\frac{1}{6}$}{$\frac{1}{4}$}{$\frac{2}{3}$}{$\frac{1}{3}$}
	\loigiai{
		Tổng số quả cầu trong hộp là 4.
		Ký hiệu: Đỏ (Đ), Xanh (X), Vàng 1 (V1), Vàng 2 (V2).
		Số cách chọn ngẫu nhiên đồng thời 2 quả cầu từ 4 quả cầu là $C_4^2 = \frac{4 \cdot 3}{2 \cdot 1} = 6$.
		Vậy số phần tử của không gian mẫu là $n(\Omega) = 6$.
		Các cặp có thể chọn là: (Đ;X), (Đ;V1), (Đ;V2), (X;V1), (X;V2), (V1;V2).
		
		Gọi A là biến cố “Chọn được 1 quả cầu đỏ và 1 quả cầu vàng”.
		Để chọn được 1 quả cầu đỏ, có 1 cách chọn (Đ).
		Để chọn được 1 quả cầu vàng, có 2 cách chọn (V1 hoặc V2).
		Số kết quả thuận lợi cho biến cố A là $n(A) = 1 \times 2 = 2$.
		Các cặp thuận lợi là: (Đ;V1), (Đ;V2).
		
		Xác suất của biến cố A là: $P(A) = \frac{n(A)}{n(\Omega)} = \frac{2}{6} = \frac{1}{3}$.
		\\ \textbf{Đáp Án D}
	}
\end{ex}

\newpage
\textbf{II. PHẦN TỰ LUẬN}
\begin{ex}
	Rút gọn biểu thức sau: $B=\left(2+\frac{x+3\sqrt{x}}{\sqrt{x}+3}\right)\left(2+\frac{x-3\sqrt{x}}{3-\sqrt{x}}\right)$, với $x \ge 0; x \neq 9$.
	\loigiai{
		Điều kiện: $x \ge 0; x \neq 9$.
		Ta có:
		\begin{align*} B &= \left(2+\frac{\sqrt{x}(\sqrt{x}+3)}{\sqrt{x}+3}\right)\left(2+\frac{\sqrt{x}(\sqrt{x}-3)}{-(\sqrt{x}-3)}\right) \\ &= (2+\sqrt{x})(2-\sqrt{x}) \quad (\text{với } x \neq 9 \text{ nên } \sqrt{x}-3 \neq 0) \end{align*}
		Hoặc, nếu ta rút gọn từng phần trong ngoặc trước:
		\begin{align*} \frac{x+3\sqrt{x}}{\sqrt{x}+3} &= \frac{\sqrt{x}(\sqrt{x}+3)}{\sqrt{x}+3} = \sqrt{x} \\ \frac{x-3\sqrt{x}}{3-\sqrt{x}} &= \frac{\sqrt{x}(\sqrt{x}-3)}{-(\sqrt{x}-3)} = -\sqrt{x} \quad (\text{do } x \neq 9 \text{ nên } \sqrt{x}-3 \neq 0) \end{align*}
		Khi đó:
		\begin{align*} B &= (2+\sqrt{x})(2+(-\sqrt{x})) \\ &= (2+\sqrt{x})(2-\sqrt{x}) \\ &= 2^2 - (\sqrt{x})^2 \\ &= 4-x \end{align*}
		Vậy $B = 4-x$.
	}
\end{ex}

\begin{ex}
	Giải hệ phương trình $\begin{cases} x+2y=1 \\ 3x+y=8 \end{cases}$.
	\loigiai{
		Ta có hệ phương trình:
		$$ \begin{cases} x+2y=1 \quad &(1) \\ 3x+y=8 \quad &(2) \end{cases} $$
		Từ phương trình (2), suy ra $y = 8-3x$.
		Thế $y = 8-3x$ vào phương trình (1):
		\begin{align*} x + 2(8-3x) &= 1 \\ x + 16 - 6x &= 1 \\ -5x &= 1 - 16 \\ -5x &= -15 \\ x &= \frac{-15}{-5} \\ x &= 3 \end{align*}
		Thay $x=3$ vào $y = 8-3x$:
		\begin{align*} y &= 8 - 3(3) \\ &= 8 - 9 \\ &= -1 \end{align*}
		Vậy nghiệm của hệ phương trình là $(x;y) = (3;-1)$.
	}
\end{ex}

\begin{ex}
	Cho phương trình $x^2-5x+6=0$ không giải phương trình, hãy tính giá trị của biểu thức $T=x_1^2x_2+x_1x_2^2$.
	\loigiai{
		Phương trình $x^2-5x+6=0$ có $a=1, b=-5, c=6$.
		$\Delta = b^2-4ac = (-5)^2 - 4(1)(6) = 25 - 24 = 1 > 0$.
		Phương trình có hai nghiệm phân biệt $x_1, x_2$.
		Theo hệ thức Vi-ét, ta có:
		$$ \begin{cases} x_1 + x_2 = -\frac{b}{a} = -\frac{-5}{1} = 5 \\ x_1 x_2 = \frac{c}{a} = \frac{6}{1} = 6 \end{cases} $$
		Ta biến đổi biểu thức T:
		\begin{align*} T &= x_1^2x_2+x_1x_2^2 \\ &= x_1x_2(x_1+x_2) \end{align*}
		Thay giá trị từ hệ thức Vi-ét vào biểu thức T:
		\begin{align*} T &= (6)(5) \\ &= 30 \end{align*}
		Vậy $T=30$.
	}
\end{ex}

\begin{ex}
	Một đội xe nhận vận chuyển 96 tấn hàng. Nhưng khi sắp khởi hành có thêm 3 xe nữa, nên mỗi xe chở ít hơn lúc đầu 1,6 tấn hàng. Hỏi lúc đầu đội xe có bao nhiêu chiếc?
	\loigiai{
		Gọi số xe lúc đầu của đội là $x$ (chiếc).
		Điều kiện: $x \in \mathbb{N}^*$.
		
		Lúc đầu, mỗi xe chở số tấn hàng là: $\frac{96}{x}$ (tấn).
		Lúc sau, số xe của đội là: $x+3$ (chiếc).
		Lúc sau, mỗi xe chở số tấn hàng là: $\frac{96}{x+3}$ (tấn).
		
		Vì lúc sau mỗi xe chở ít hơn lúc đầu 1,6 tấn hàng, ta có phương trình:
		$$ \frac{96}{x} - \frac{96}{x+3} = 1,6 $$
		Đổi $1,6 = \frac{16}{10} = \frac{8}{5}$.
		$$ \frac{96}{x} - \frac{96}{x+3} = \frac{8}{5} $$
		Quy đồng mẫu số và giải phương trình:
		\begin{align*}
			96 \cdot 5 (x+3) - 96 \cdot 5 x &= 8x(x+3) \\
			480(x+3) - 480x &= 8x(x+3) \\
			480x + 1440 - 480x &= 8x^2 + 24x \\
			1440 &= 8x^2 + 24x \\
			8x^2 + 24x - 1440 &= 0
		\end{align*}
		Chia cả hai vế cho 8:
		$$ x^2 + 3x - 180 = 0 $$
		Giải phương trình bậc hai $x^2 + 3x - 180 = 0$:
		Ta có $\Delta = b^2 - 4ac = 3^2 - 4(1)(-180) = 9 + 720 = 729$.
		$\sqrt{\Delta} = \sqrt{729} = 27$.
		Phương trình có hai nghiệm phân biệt:
		$x_1 = \frac{-b + \sqrt{\Delta}}{2a} = \frac{-3 + 27}{2 \cdot 1} = \frac{24}{2} = 12$.
		$x_2 = \frac{-b - \sqrt{\Delta}}{2a} = \frac{-3 - 27}{2 \cdot 1} = \frac{-30}{2} = -15$.
		
		Đối chiếu với điều kiện $x \in \mathbb{N}^*$:
		Nghiệm $x_1 = 12$ (thỏa mãn).
		Nghiệm $x_2 = -15$ (không thỏa mãn).
		
		Vậy, lúc đầu đội xe có 12 chiếc.
	}
\end{ex}
\begin{ex}
	Một hộp có 30 chiếc thẻ cùng loại, mỗi thẻ được ghi một trong các số 2, 4, 6, ..., 60, hai thẻ khác nhau ghi hai số khác nhau. Xét phép thử “Rút ngẫu nhiên một thẻ trong hộp”. Tính xác suất của biến cố sau: A: “Số xuất hiện trên thẻ được rút ra lớn hơn 12 và là ước của 60”.
	\loigiai{
		Các số ghi trên thẻ là các số chẵn từ 2 đến 60.
		Không gian mẫu $\Omega = \{2; 4; 6; 8; 10; 12; 14; 16; 18; 20; 22; 24; 26; 28; 30; 32; 34; 36; 38; 40; 42; 44; 46; 48; 50; 52; 54; 56; 58; 60\}$.
		Số phần tử của không gian mẫu là $n(\Omega) = 30$.
		
		Biến cố A: “Số xuất hiện trên thẻ được rút ra lớn hơn 12 và là ước của 60”. các kết quả thuận lợi cho biến cố A là:
		$A = \{20; 30; 60\}$.
		Số phần tử của biến cố A là $n(A) = 3$.
		
		Xác suất của biến cố A là:
		$P(A) = \frac{n(A)}{n(\Omega)} = \frac{3}{30} = \frac{1}{10}$.
	}
\end{ex}
\begin{ex}
	Cho đường tròn (O) có đường kính AB và điểm C thuộc đường tròn đó (C khác A, B). Lấy điểm D thuộc dây BC (D khác B, C). Tia AD cắt cung nhỏ BC tại điểm E, tia AC cắt tia BE tại điểm F.
	\begin{enumerate}
		\item[a)] Chứng minh rằng 4 điểm F, C, D, E thuộc một đường tròn.
		\item[b)] Chứng minh rằng $DA \cdot DE = DB \cdot DC$. Gọi I là tâm đường tròn ngoại tiếp tứ giác FCDE, chứng minh rằng IC là tiếp tuyến của đường tròn (O).
	\end{enumerate}
\end{ex}

\begin{ex}
	Cho các số thực dương a, b, c thỏa mãn $20a + 15b + 12c = abc$.
	Tìm giá trị lớn nhất của biểu thức $P = \frac{3}{\sqrt{a^2+9}} + \frac{4}{\sqrt{b^2+16}} + \frac{5}{\sqrt{c^2+25}}$.
\end{ex}

\setcounterref{ex}{0}
